
\label{sec:problema}


\section{Planteamiento del problema}

Las estimaciones oficiales del DANE se desagregan, por lo general, a nivel departamental o para las principales áreas metropolitanas. Esta escala analítica diluye las heterogeneidades locales y oculta los contrastes socioeconómicos existentes entre municipios \cite{DANE_PobrezaMultidimensional_2023}. La ausencia de datos confiables a escala municipal debilita la focalización de programas sociales, la distribución eficiente de recursos y la formulación de políticas adaptadas a las realidades de las regiones históricamente rezagadas. La raíz del problema es metodológica: las encuestas de hogares empleadas para estas mediciones como la Gran Encuesta Integrada de Hogares (GEIH) o la Encuesta Nacional de Calidad de Vida (ECV), poseen diseños muestrales representativos a nivel departamental, pero carecen de la potencia estadística necesaria para realizar inferencias válidas en el ámbito municipal \cite{DANE_ENCV_2022}.

A esta restricción espacial se suma una limitación temporal. Colombia carece de mediciones oficiales directas y frecuentes del IPM municipal, lo que impide monitorear oportunamente la evolución de las privaciones en unidades territoriales pequeñas. Las estadísticas oficiales con mayor desagregación geográfica provienen de los censos de población y vivienda debido a su cobertura exhaustiva. Sin embargo, al ser operaciones periódicas y no anuales, los censos no permiten un seguimiento continuo. Este desfase intercensal genera un vacío de información que restringe la actualización del IPM municipal y reduce la capacidad institucional para detectar cambios recientes en las condiciones de vida de la población \cite{DANE_PobrezaMultidimensional_2023, DANE_ENCV_2022}.

Adicionalmente, existe una debilidad metodológica que se origina en la carencia de representatividad muestral de las encuestas oficiales como la GEIH y la ECV, las cuales están diseñadas para producir estimaciones departamentales o regionales, pero no locales \cite{DANE_ENCV_2022}. Esta limitación impide identificar desigualdades territoriales específicas y restringe la capacidad institucional para focalizar políticas públicas efectivas en municipios con alta vulnerabilidad socioeconómica \cite{CEPAL_PanoramaSocial_2023}.

En un entorno de alta desigualdad y heterogeneidad territorial como el colombiano, los promedios departamentales suelen enmascarar brechas críticas; dentro de un mismo departamento coexisten municipios con estándares de bienestar urbanos junto a otros con niveles de pobreza propios de las zonas rurales más vulnerables \cite{CEPAL_PanoramaSocial_2023}.

Si la falta de estimaciones municipales precisas se mantiene, las implicaciones serán concretas y de alto impacto para la gestión pública: i) la asignación de recursos continuará guiándose por promedios departamentales, lo que puede ocultar diferencias relevantes entre municipios y dejar sin focalización adecuada a territorios altamente vulnerables; ii) programas sociales y proyectos de inversión podrían seguir implementándose en zonas no prioritarias, reduciendo la eficacia y la eficiencia del gasto público; y iii) la ausencia de información granular a nivel municipal limita la capacidad institucional de monitorear intervenciones, evaluar resultados y ajustar políticas basadas en evidencia. En consecuencia, la falta de indicadores confiables a esta escala perpetúa decisiones menos informadas y debilita la capacidad del Estado para responder de manera diferenciada a los desafíos territoriales.
